% !TEX program = lualatex
% !TEX TS-program = lualatex
% main.tex: your document. Build and check it with
%   make -f Makefile.a11y check
% example.tex has every kind of content, tagged and explained. Find the block
% you need by its banner line (search for "%----") and copy it here.

%------------------------------------------------------------------------------
% METADATA AND CLASS
%------------------------------------------------------------------------------
% Keep this block above \documentclass: it turns tagging on and names the
% standard. Meets: ISO 14289-2:2024 8.11 (metadata).
\DocumentMetadata{
  lang        = en-US,
  pdfstandard = ua-2,
  tagging     = on,
  check-tagging-status,
}
% Use [report] for chapters. [book] also works, but book has no abstract, so
% delete the abstract below. Clemson's text concept prefers left-aligned text,
% so align=ragged is the default; align=justified overrides it. README.md lists
% the other options. A misspelled key=value option stops the build.
\documentclass{clemsona11y}

%------------------------------------------------------------------------------
% PACKAGES
%------------------------------------------------------------------------------
% The class already loads amsmath, amsthm, graphicx, hyperref, enotez, fontspec
% and unicode-math. Add only what you use; README.md ("Packages by field")
% says which packages tag correctly and which to avoid.
\usepackage{booktabs}  % \toprule, \midrule and \bottomrule in tables

%------------------------------------------------------------------------------
% PICTURES
%------------------------------------------------------------------------------
% Put every figure and image in the resources/ folder next to this file.
% \includegraphics[alt={...}]{name} finds it there by file name. LaTeX looks
% there only for pictures: references.bib and any \input file stay next to
% main.tex. The example's pictures are there too; example.tex still needs them.

%------------------------------------------------------------------------------
% TITLE AND ABSTRACT
%------------------------------------------------------------------------------
% pdftitle and pdfauthor set the PDF metadata. The doubled braces keep a comma
% from splitting the title. Separate the names in pdfauthor with commas.
\title[pdftitle={{Title of the document}}]{Title of the document}
\author[pdfauthor={First Author, Second Author}]{First Author and Second Author}
\date{\today}

\begin{document}
\maketitle

\begin{abstract}
One paragraph that says what the document is about.
\end{abstract}

%------------------------------------------------------------------------------
% CONTENTS AND LISTS
%------------------------------------------------------------------------------
% Uncomment the ones the document needs.
% \tableofcontents
% \listoffigures
% \listoftables

%------------------------------------------------------------------------------
% BODY
%------------------------------------------------------------------------------
\section{Introduction}
Write here. Use \verb|\section|, \verb|\subsection| and \verb|\subsubsection| in order
and never skip a level. Give every picture alt text and declare the header cells of every
table; \texttt{example.tex} shows both.

%------------------------------------------------------------------------------
% ENDNOTES AND BIBLIOGRAPHY
%------------------------------------------------------------------------------
% \printendnotes prints the endnotes written with \endnote. Uncomment the four
% bibliography lines once something is cited; with no \cite they print an
% empty References heading. Entries go in references.bib. make runs BibTeX.
% VS Code runs it once you add "% !BIB program = bibtex" and
% "% !BIB TS-program = bibtex" under line 2. In TeXShop, run Typeset, BibTeX,
% Typeset, Typeset by hand.
\printendnotes
% \phantomsection
% \addcontentsline{toc}{section}{\refname}  % report/book: {chapter}{\bibname}
% \bibliographystyle{plain}
% \bibliography{references}
\end{document}
